% ============================================================
%  part3/introduction.tex  —  Introduction to Part III
% ============================================================

\chapter*{Introduction to Part III}
\addcontentsline{toc}{chapter}{Introduction to Part III}
\markboth{Introduction to Part III}{Introduction to Part III}

Part~III is where the book shifts from reconstruction theory
to dynamics.
The reader has learned that observations collapse equivalence
classes.
Now they learn that physical systems themselves naturally evolve
toward lower-dimensional descriptions through coarsening.

\begin{tcolorbox}[
  title={\textbf{Established So Far}},
  colback=boxgreen,colframe=green!50!black,
  fonttitle=\small\sffamily
]
\begin{itemize}[leftmargin=1.5em,itemsep=2pt]
  \item Observation is projection: $\pi : X \to M$ with
    $\ker(\pi) \neq 0$.
  \item Reconstruction is imperfect whenever $F$ is
    non-injective: $E(x) > 0$.
  \item CLIO minimises reconstruction error subject to
    admissibility constraints.
  \item Gradient flows satisfy $\frac{d\mathcal{F}}{dt} \leq 0$.
  \item Entropy production is nonnegative.
\end{itemize}
\end{tcolorbox}

\bigskip

\begin{tcolorbox}[
  title={\textbf{New Ideas Introduced in This Part}},
  colback=boxblue,colframe=blue!40!black,
  fonttitle=\small\sffamily
]
\begin{itemize}[leftmargin=1.5em,itemsep=2pt]
  \item \textbf{Symmetry breaking}: systems with symmetric
    Hamiltonians that choose asymmetric ground states.
  \item \textbf{Domain formation}: spatial regions of distinct
    phases separated by interfaces.
  \item \textbf{Curvature-driven coarsening}: $V_n \propto -\kappa$
    and the resulting law $\ell(t) \sim t^{1/z}$.
  \item \textbf{Allen--Cahn dynamics}: $z = 2$, non-conserved.
  \item \textbf{Cahn--Hilliard dynamics}: $z = 3$, conserved.
  \item \textbf{Dynamic scaling}: $C(r,t) = f(r/\ell(t))$.
  \item \textbf{Topological defects}: vortices, walls, strings,
    nodes classified by homotopy groups.
  \item \textbf{Universality}: microscopic details are irrelevant;
    only symmetry and conservation determine $z$.
  \item \textbf{Cross-domain examples}: river basins, foams,
    ecological networks, neural synchronisation, cosmic web.
\end{itemize}
\end{tcolorbox}

\bigskip

\begin{tcolorbox}[
  title={\textbf{Dependencies}},
  colback=boxgray,colframe=gray!50!black,
  fonttitle=\small\sffamily
]
Part~III requires: basic calculus, the variational derivative
from Chapter~\ref{ch:variational}, and the gradient-flow
dissipation result from Chapter~\ref{ch:fields-flows}.
No prior knowledge of condensed matter physics is assumed.
\end{tcolorbox}

\bigskip

\begin{tcolorbox}[
  title={\textbf{What Remains Open After This Part}},
  colback=boxorange,colframe=orange!60!black,
  fonttitle=\small\sffamily
]
\begin{itemize}[leftmargin=1.5em,itemsep=2pt]
  \item The RSVP coarsening exponent $z_{\mathrm{RSVP}}$
    is conjectured but not yet derived. (Chapter~\ref{ch:open-problems})
  \item The relationship between topological defects and
    sheaf cohomology obstructions is established conceptually
    here but formalised only in
    Chapter~\ref{ch:sheaf}.
  \item Whether the cosmic web belongs to the same
    universality class as the cross-domain examples is
    the central empirical question of Part~VII\@.
\end{itemize}

The central thesis of Part~III:
\principle{
  Structure emerges through the progressive elimination\\[4pt]
  of microscopically inadmissible distinctions.
}
This is the dynamical counterpart of Part~I's observation
that representation requires forgetting.
\end{tcolorbox}
